% sec-00: abstract, executive summary, reader's guide, and the two indexes.
% DRAFT stage.  Tables 1 and 2 are written in full here, because figure and
% table numbering is fixed at this stage and never moves again.
\section*{Abstract}
\phantomsection\addcontentsline{toc}{section}{Abstract}

\draftinstr{Write 200 to 240 words. Open with the conversion: the author has
filed ten funding applications as an independent scientist, and one of the ten,
application 05, is addressed to the only mechanism in the set built for a
company. State the three numbers that govern the paper, \$306,000, \$1,300,000
and \$3,500,000, and the two derived ones, \$1,396,000 of direct work inside the
award and \$2,104,000 outside it. Name the clause: not the report's Pioneer
language and not its NOVEL PERFORMERS chapter, but the SBIR clause, which
appears in Chapters I, III and IV. Close on what is being offered, which is a
costed twelve-milestone schedule with an evidence artifact and a stop condition
on every row. Source: \appfile{funding/pdac-funding-applications/final-apply/sections/sec-00-front.tex}
for the abstract form, and
\appfile{funding/science-golden-age/chunk-03-chapter-two-revitalizing-science-and-technology-enterprise.md}
for the clause.}

\section*{Executive Summary}
\phantomsection\addcontentsline{toc}{section}{Executive Summary}

\draftinstr{Six short paragraphs, one per outline item A to E plus one on what
the conversion costs. Each must carry at least one number. Do not restate the
abstract. Sources:
\appfile{funding/pdac-funding-applications/final-apply/sections/sec-08-budget-and-leverage.tex}
for the budget frame,
\appfile{funding/pdac-funding-applications/applications/app-05-nih-sbir-seed/}
for the award split, and
\appfile{funding/science-golden-age/chunk-08-annex-fiscal-year-2028-research-and-development-budget-priorities.md}
for the 3 to 1 leverage target.}

\section*{How to Read This Document}
\phantomsection\addcontentsline{toc}{section}{How to Read This Document}

\draftinstr{One short paragraph plus a four-row table mapping four readers to
the sections each should read first: an SBIR program director, an ARPA-H
programme manager, a private investor, and a UC San Diego intake committee.
Source: \appfile{funding/pdac-funding-applications/final-apply/sections/sec-02-ten-applications.tex}
for the reader-mapping form.}

\begin{apptable}
\begin{tabularx}{\textwidth}{>{\raggedright\arraybackslash}p{0.7cm} >{\raggedright\arraybackslash}p{1.1cm} >{\raggedright\arraybackslash}p{2.5cm} >{\raggedright\arraybackslash}p{3.3cm} >{\raggedright\arraybackslash}X}
\headrow \thead{Fig} & \thead{\S} & \thead{Platform} & \thead{Source directory} & \thead{The question it answers}\\
\midrule
1 & 1 & Mermaid & \appfile{mermaid/} & Which clause does a company with no indirect-cost base qualify under? \\
2 & 1 & D2 & \appfile{d2/} & What kind of institution is a 2.6 FTE firm, on the report's own axes? \\
3 & 1 & Graphviz & \appfile{graphviz/} & What does the same direct work cost under three overhead regimes? \\
4 & 2 & Diagrams & \appfile{diagrams-python/} & What is the shape of what the company holds today? \\
5 & 2 & D2 & \appfile{d2/} & What is owned, licensed, contracted, and absent, with terms? \\
6 & 2 & PlantUML & \appfile{plantuml/} & What may the sponsor do alone, and what needs an institution? \\
7 & 3 & Mermaid & \appfile{mermaid/} & What must hold at month nine for Phase II to begin? \\
8 & 3 & D2 & \appfile{d2/} & What does the same programme cost under two mechanisms? \\
9 & 3 & Graphviz & \appfile{graphviz/} & Which work packages does the money reach, and which not? \\
10 & 4 & PlantUML & \appfile{plantuml/} & Which capital positions may be held while a participant is on study? \\
11 & 4 & D2 & \appfile{d2/} & Where does every dollar come from, and what separates the tiers? \\
12 & 4 & Mermaid & \appfile{mermaid/} & In what order is a private round executed during enrollment? \\
13 & 5 & Mermaid & \appfile{mermaid/} & When does each milestone open, close, and produce its artifact? \\
14 & 5 & Graphviz & \appfile{graphviz/} & What has to fail, and in what combination, to stop the programme? \\
15 & 5 & PlantUML & \appfile{plantuml/} & What do three parties do at once to close one milestone? \\
16 & 6 & D2 & \appfile{d2/} & Which published quantities can a funder check before committing? \\
17 & 6 & Graphviz & \appfile{graphviz/} & Where does each published number land in a filing? \\
18 & 7 & Diagrams & \appfile{diagrams-python/} & Who is employed, who is contracted, who is contributed? \\
19 & 8 & Mermaid & \appfile{mermaid/} & What did the July and August 2026 contacts actually unlock? \\
20 & 10 & Diagrams & \appfile{diagrams-python/} & What survives if the programme stops, and who holds it? \\
\end{tabularx}
\end{apptable}
\tabcap{The twenty figures, their section, their platform, the directory that
carries each specification, and the one question each answers. Five mermaid,
three plantuml, five d2, three diagrams, four graphviz, by purpose not quota.}

\begin{apptable}
\begin{tabularx}{\textwidth}{>{\raggedright\arraybackslash}p{0.7cm} >{\raggedright\arraybackslash}p{1.1cm} >{\raggedright\arraybackslash}p{4.6cm} >{\raggedright\arraybackslash}X}
\headrow \thead{Tab} & \thead{\S} & \thead{Contents} & \thead{The question it answers}\\
\midrule
1 & 0 & Figure index, twenty rows & Where is each figure and what does it answer? \\
2 & 0 & Table index, nineteen rows & Where is each table and what does it answer? \\
3 & 0 & Four readers, four entry points & Which section should this reader open first? \\
4 & 1 & Three clauses against four eligibility tests & Which clause survives the filter, and why? \\
5 & 1 & Three overhead regimes on one direct base & What is the premium for routing through an institution? \\
6 & 2 & The owned register, thirteen rows with DOIs & What exists today and where is it deposited? \\
7 & 2 & The absent register, seven rows with signatories & What is missing and who alone can supply it? \\
8 & 3 & The four-layer budget under two prices & What does the SBIR route reach, and what does it not? \\
9 & 3 & Phase I line items, eight rows to \$306,000 & What does nine months of Phase I actually buy? \\
10 & 3 & Phase II line items, ten rows to \$1,300,000 & What does twenty-four months of Phase II buy? \\
11 & 4 & The \S54.2 and \S54.4 disclosure triggers & What must be disclosed, at what threshold, on which form? \\
12 & 4 & The capitalization table across five stages & Who owns what, and does the SBIR test still pass? \\
13 & 5 & The twelve milestones, costed and gated & What is the evidence, the cost, and the stop condition? \\
14 & 5 & Milestone artifacts against funder mechanism & Which mechanism pays for which milestone? \\
15 & 6 & Six checkable quantities with limitations & What can a reviewer verify against a published source? \\
16 & 6 & Cost and schedule against industry benchmarks & What did the prior work cost, and what does it cost elsewhere? \\
17 & 7 & The 2.6 FTE by employment status & Who is on the payroll and who is not? \\
18 & 9 & Risks, likelihood, and the mitigation held & What could go wrong and what is already in place? \\
19 & 11 & Abbreviations, two column pairs & What does this term mean in this paper? \\
\end{tabularx}
\end{apptable}
\tabcap{The nineteen tables, their section, their contents, and the one question
each answers. Every table is set to the full body measure with ragged-right
fixed columns, so no cell shows a stretched interword gap.}
